\documentclass[11pt]{article}

\usepackage[margin=1in]{geometry}
\usepackage{microtype}
\usepackage{amsmath}
\usepackage{graphicx}
\usepackage{booktabs}
\usepackage[hidelinks]{hyperref}

\graphicspath{{figures/}}

\newcommand{\natex}{\texttt{natex}}

\title{Teaching to the Test, Undetected: No Release-Localized Contamination
Premium on Public versus Held-Out Benchmarks}
\author{Hauke Hillebrandt\thanks{University College London. Email:
\texttt{ucjthhi@ucl.ac.uk}. This paper and the underlying \natex{} software were
prepared with substantial assistance from Anthropic's Claude models; the author
reviewed the analyses and text and is responsible for all remaining errors.}}
\date{July 2026 \\ \small analysis pass of record: 2026-07, seed 0 (wild-cluster
bootstrap only; randomization inference is exhaustive)}

\begin{document}

\maketitle

\begin{abstract}
\noindent In Epoch AI's cross-benchmark capability panel, models released
after a benchmark's publication score above their general-capability (ECI)
expectation on \emph{public} benchmarks relative to four \emph{held-out}
benchmarks --- but the premium does not survive cluster-honest inference, and
the design's own placebo battery shows it is not localized at benchmark
release. In 896 model~$\times$~benchmark cells spanning the 19 benchmarks
that straddle their own release date, the post~$\times$~public level contrast
is $\tau=+0.127$~z (CR1 SE $0.079$); every level specification is positive,
and late models underperform their ECI expectation on held-out benchmarks by
$-0.19$~z --- exactly the sign pattern a teaching-to-the-test story predicts.
But with only 4 held-out benchmark clusters the inference of record is a
wild-cluster bootstrap ($p=0.145$) and exact randomization inference over all
$\binom{19}{4}=3{,}876$ held-out labelings (one-sided $p=0.213$), neither of
which rejects. The premium collapses under IRT-difficulty controls ($+0.04$)
and model fixed effects ($+0.03$) --- held-out status is confounded with
benchmark difficulty --- and it fails the placebo-cutoff localization test:
the contrast at a $-0.5$-year placebo cutoff ($+0.267$) \emph{exceeds} the
true-cutoff estimate, the signature of a smooth group trend difference in
event time rather than a step at release. Ability-proxy variants that avoid
ECI's mechanical attenuation (the index is trained on the public benchmarks)
triple the estimate ($+0.35$ to $+0.39$, wild-cluster $p=0.047$--$0.082$) but
inherit the same localization failure. A difference-in-kinks secondary is a
correctly identified artifact (its placebo cutoff beats every true-cutoff
estimate), and the automated \natex{} SuDDDS run recovered the declared
contrast but honestly \emph{refused} a placebo $p$ (no in-space placebos with
a single binary group dimension); the manual bootstrap/randomization battery
fills exactly that gap and does not reject. The verdict is a suggestive-sign
null: no release-localized contamination premium is detectable in this panel.
\end{abstract}

\section{Introduction}\label{sec:intro}

Train--test contamination --- benchmark answers leaking into the training
corpora of the models being benchmarked --- is the standing objection to
public AI leaderboards. \cite{oren2023} show that test-set membership can be
proven from a model's preference for the canonical ordering of an
exchangeable benchmark; \cite{roberts2023} exploit the natural experiment of
GPT training cutoffs to show that performance on code benchmarks is
systematically related to whether the problems predate the cutoff;
\cite{zhang2024} commission a fresh clone of GSM8k and find accuracy drops of
up to 13\% for model families that apparently overfit the public original.
\cite{dominguez2024} generalize the worry beyond verbatim leakage:
\emph{training on the test task} --- benchmark-shaped data and formats
entering training pipelines --- confounds model comparisons even when no
test item was ever seen.

The benchmark-building community's institutional response has been to hold
answers back. FrontierMath keeps private problem tiers scored only via Epoch
AI \cite{glazer2024}; Humanity's Last Exam withholds a private split against
overfitting \cite{phan2025}; ARC-AGI-2 maintains a private evaluation set
\cite{chollet2025}; LiveBench rotates fresh questions monthly precisely to
stay ahead of training corpora \cite{white2024}. This design pattern implies
a cross-benchmark test of contamination in the wild. If public-benchmark
scores are inflated by leakage or test-task training, then models released
\emph{after} a public benchmark's publication should overperform their
general capability on it, and the premium should be absent on held-out
benchmarks whose answers cannot be downloaded --- a
difference-in-differences in event time around each benchmark's release
date, with held-out benchmarks as the comparison group.

This paper runs that test on Epoch AI's benchmarking panel
\cite{epochai,ho2025}, treating benchmark release as a dated event inside a
model~$\times$~benchmark panel and demanding that any premium be
\emph{localized} at the event: estimated with cluster-honest inference (only
4 of 19 benchmark clusters are held out), stress-tested against difficulty
controls, and benchmarked against placebo release dates, with the automated
\natex{} toolkit's discovery-and-refusal machinery \cite{natex} reported as
run. The sign pattern survives everything; the localized premium survives
nothing.

\section{Data}\label{sec:data}

Three public Epoch AI files are joined \cite{epochai}: the ECI benchmark
score table (2{,}010 model~$\times$~benchmark performance cells with
benchmark release dates), the ECI model index (211 models with release dates
and ECI capability scores \cite{ho2025}), and the IRT benchmark parameter
table (52 benchmarks with estimated difficulty \texttt{edi} and
discrimination). Three release dates missing from the source are hand-filled
as declared inputs: HLE 2025-01-23, ARC-AGI-2 2025-03-24, and
FrontierMath-Tier-4-Private 2025-07-01. The held-out classification is
likewise declared, by whether answers are publicly downloadable:
FrontierMath-2025-02-28-Private, FrontierMath-Tier-4-2025-07-01-Private,
HLE, and ARC-AGI-2 (4 benchmarks); all other benchmarks are public.

A cell is a model~$\times$~benchmark pair. The running variable is event
time, $x = $ model release date $-$ benchmark release date, in years. The
identifying sample keeps the 19 benchmarks with cells on both sides of
$x=0$: $n=896$ cells across 150 models, split 227 pre-release / 668
post-release (one cell at exactly $x=0$ is coded pre). Outcomes are built
within benchmark: $z$ is the within-benchmark z-score of performance, and
the primary outcome $r$ residualizes $z$ linearly on the model's ECI score,
so $r$ measures performance relative to what the model's general capability
predicts on that benchmark. Variants used below: a logit-performance
residual (48 zero and 3 one scores clipped at $[0.01,0.99]$), and two
ability proxies that avoid ECI itself --- the model's mean $z$ across all
\emph{other} benchmarks (leave-one-out) and across held-out benchmarks only.
The panel builder asserts bit-identical reproduction of the scout-pass $r$
before writing any file.

\section{Design and Methods}\label{sec:methods}

\paragraph{Level contrast.} The primary estimand is the post-release premium
of public over held-out benchmarks in $r$, from the cell-level regression
\begin{equation*}
r_{mb} = \alpha + \beta x_{mb} + \gamma\,\mathrm{post}_{mb}
+ \tau\,(\mathrm{post}_{mb} \times \mathrm{public}_b) + \mu_b + u_{mb},
\end{equation*}
with benchmark fixed effects $\mu_b$ and a common linear event-time trend;
$\gamma$ is the held-out post-release baseline and $\tau$ the contamination
premium. Errors are CR1 cluster-robust by benchmark with a $t(18)$
reference. With 19 clusters of which only 4 are held out, CR1 alone is not
trustworthy \cite{cameron2015}, so the inference of record is (i) a
null-imposed Rademacher wild-cluster bootstrap ($B=9{,}999$, seed 0) and
(ii) \emph{exact} randomization inference over all $\binom{19}{4}=3{,}876$
assignments of the held-out label to 4 of the 19 benchmarks,
FWL-accelerated (the observed labeling's coefficient is ranked in the full
permutation distribution; no sampling anywhere).

\paragraph{Localization.} A contamination premium must be a step at
benchmark release, not a generic divergence of the two groups in event
time. The identical regression is re-estimated with the cutoff moved to
$x_0 \in \{-1.0, -0.5, +0.5, +1.0\}$ years; a true release-localized effect
should peak at the true cutoff. Side composition (straddling benchmarks and
held-out benchmarks per placebo) is reported because extreme placebo cutoffs
thin one side.

\paragraph{Robustness.} Seven re-estimates: the logit-residual outcome;
post~$\times$~IRT-difficulty and post~$\times$~discrimination controls
(standardized \texttt{edi} and slope); a two-way model~$+$~benchmark
fixed-effects specification on $z$; the two ability-proxy outcomes;
dropping saturated cells (performance $>0.95$); and model-release-year fixed
effects, which isolate cross-benchmark stagger from calendar time. A
per-benchmark jackknife re-estimates $\tau$ dropping one cluster at a time.

\paragraph{Automated discovery (SuDDDS).} The \natex{} \texttt{discover}
pipeline \cite{natex} ran on the benchmark $\times$ half-year-bin panel
(94 cells; \texttt{--design did}, declared treatment path
$\theta = \mathrm{public} \times \mathrm{post}$, seed 0). Its subset scan,
anticipation gate, and
placebo machinery either validate the declared event or refuse; refusals are
reported as run, per the toolkit's honest-inference conventions.

\paragraph{Difference-in-kinks secondary.} As a slope-based cross-check, a
group difference-in-kinks in the tradition of \cite{bockerman2025} ---
\natex{} \texttt{kink} with the held-out group as the comparison dimension,
triangular kernel, CR1 by benchmark --- estimates the change in the
event-time \emph{slope} contrast at $x=0$, over bandwidths $0.5$--$2.0$
years, with a placebo cutoff at $x_0=+0.5$.

\section{Results}\label{sec:results}

\begin{figure}[t]
\centering
\includegraphics[width=\textwidth]{fig1}
\caption{(a) Mean ECI-residualized within-benchmark z-score $r$ by 0.5-year
event-time bin (whiskers $\pm1$ SE; bins with $\geq 3$ cells), public versus
held-out benchmarks, vertical rule at benchmark release. In the first year
after release, held-out cells fall below their capability expectation while
public cells sit above it --- but the raw profiles show no step at $x=0$,
and the two groups already differ before release. (b) The
post~$\times$~public contrast $\tau$ across specifications
(filled circles; CR1 95\% CIs, $t(18)$ critical value; wild-cluster
bootstrap $p$ annotated where computed) and at placebo event-time cutoffs
(open circles). Every specification is positive, but the $-0.5$-year placebo
cutoff exceeds the true-cutoff estimate --- the localization failure that
grades this a null.}
\label{fig:fig1}
\end{figure}

\paragraph{A positive premium with a contamination-consistent sign pattern.}
The primary contrast is $\tau=+0.1267$~z (CR1 SE $0.0786$, $t=1.61$,
$p_{t(18)}=0.124$) on $n=896$ cells in 19 benchmark clusters. The held-out
post-release baseline is $\gamma=-0.1929$~z (SE $0.0526$): after a
benchmark's release, models underperform their ECI expectation on held-out
benchmarks, while public benchmarks absorb the offsetting premium ---
the two-sided pattern teaching-to-the-test predicts. The premium is not
driven by any single cluster: the per-benchmark jackknife keeps
$\tau \in [+0.085, +0.171]$, positive in all 19 leave-one-out samples.

\paragraph{Cluster-honest inference does not reject.} The wild-cluster
bootstrap gives $p=0.145$. Exact randomization inference over all $3{,}876$
held-out labelings gives one-sided $p=0.213$ (two-sided $0.384$): more than
a fifth of all possible ``held-out'' labelings produce a premium at least as
large as the observed one, and the permutation distribution's 5th/50th/95th
percentiles ($-0.224$ / $+0.028$ / $+0.200$) comfortably bracket
$+0.127$. With 4 treated-side clusters, this is the honest finite-sample
yardstick, and the observed contrast is unremarkable against it.

\paragraph{The premium is difficulty in disguise.} The held-out set ---
FrontierMath tiers, HLE, ARC-AGI-2 --- is also the hardest-benchmark set.
Adding post~$\times$~IRT-difficulty and post~$\times$~discrimination
controls collapses $\tau$ from $+0.127$ to $+0.044$ (wild-cluster $p=0.769$);
two-way model~$+$~benchmark fixed effects give $+0.026$ ($p=0.814$); the
logit outcome gives $+0.087$ ($p=0.334$). Two specifications move the other
way: the ability-proxy outcomes, which replace ECI --- itself fit
\emph{on the public benchmarks} \cite{ho2025} and therefore mechanically
absorbing part of any public-benchmark inflation into measured capability
--- with leave-one-out or held-out-only mean $z$. These triple the estimate
($+0.388$, wild-cluster $p=0.047$; $+0.354$, $p=0.082$), confirming the
attenuation channel is real, but they inherit the localization failure
below and the 4-cluster fragility. Dropping saturated cells ($+0.132$) and
release-year fixed effects ($+0.169$, $p=0.072$) leave the picture
unchanged.

\paragraph{Localization fails.} Moving the cutoff half a year \emph{before}
benchmark release yields $\tau=+0.267$ (SE $0.122$, 15 straddling
benchmarks, all 4 held-out) --- \emph{larger} than the true-cutoff
estimate. The grid declines monotonically through the true cutoff:
$+0.800$ at $-1.0$ (uninterpretable: 11 pre cells, 4 straddling
benchmarks), $+0.267$ at $-0.5$, $+0.127$ at $0$, $+0.046$ at $+0.5$,
$-0.311$ at $+1.0$. A release-localized contamination premium should peak
at release; a smooth difference in the two groups' event-time trends
produces exactly this monotone drift. The design therefore identifies the
positive contrast as a trend difference, not a step at benchmark release.

\paragraph{Automated runs.} The \natex{} SuDDDS pipeline discovered the
declared subset $\{\mathrm{group{:}\,public}\}$ at $t_0=0$ with window 2
(treatment-path scan $p=0.010$ --- mechanical given the declared $\theta$)
and estimated DD $\tau=+0.232$ (SE $0.088$) on the 94-cell binned panel,
but \emph{refused} its placebo randomization $p$: with a single binary
group dimension there are 0 usable in-space placebos ($\geq 5$ required),
and the anticipation gate failed closed (null Holm $p$-values). The manual
wild-cluster-bootstrap and exact-randomization battery above is precisely
the placebo inference the refusal requests, and it does not reject. The
difference-in-kinks secondary is an identified artifact: the slope contrast
at the true cutoff is $+0.015$ (SE $0.443$) at bandwidth $0.75$, drifting
to $-0.843$ (SE $0.218$) at bandwidth $2.0$ as saturated public
post-release profiles compress, while the $+0.5$ placebo cutoff gives
$-1.467$ (SE $0.644$) --- larger in magnitude than every true-cutoff
estimate --- so the slope design carries no evidential weight here.

\begin{table}[t]
\centering
\small
\caption{Headline estimates. Outcome $r$ is the within-benchmark z-score of
performance residualized on model ECI; $\tau$ is the post~$\times$~public
coefficient with benchmark FE and a linear event-time trend, CR1 clustered
by benchmark (19 clusters, 4 held-out). WCB: null-imposed Rademacher
wild-cluster bootstrap, $B=9{,}999$, seed 0. RI: exact randomization
inference over all $\binom{19}{4}=3{,}876$ held-out labelings.}
\label{tab:main}
\begin{tabular}{lrrrl}
\toprule
Quantity & $\tau$ & CR1 SE & $p$ & notes \\
\midrule
\multicolumn{5}{l}{\emph{Panel A: primary contrast and inference of record}}\\
post $\times$ public (primary) & $+0.127$ & $0.079$ & $0.145$ & WCB; $p_{t(18)}=0.124$ \\
\quad exact RI, one-sided      & ---      & ---     & $0.213$ & two-sided $0.384$ \\
held-out post baseline $\gamma$ & $-0.193$ & $0.053$ & --- & contamination-consistent sign \\
\midrule
\multicolumn{5}{l}{\emph{Panel B: robustness (WCB $p$ where computed)}}\\
logit score                 & $+0.087$ & $0.083$ & $0.334$ & $n=896$ \\
IRT-difficulty controls     & $+0.044$ & $0.122$ & $0.769$ & post $\times$ edi, post $\times$ slope \\
model $+$ benchmark FE      & $+0.026$ & $0.101$ & $0.814$ & outcome $z$ \\
ability: leave-one-out      & $+0.388$ & $0.149$ & $0.047$ & $n=895$ \\
ability: held-out only      & $+0.354$ & $0.147$ & $0.082$ & $n=661$ \\
drop saturated (perf $>0.95$) & $+0.132$ & $0.082$ & --- & $n=878$ \\
release-year FE             & $+0.169$ & $0.081$ & $0.072$ & $n=896$ \\
\midrule
\multicolumn{5}{l}{\emph{Panel C: placebo event-time cutoffs (straddling / held-out)}}\\
cutoff $-1.0$ yr & $+0.800$ & $0.078$ & --- & 4/1; 11 pre cells, uninterpretable \\
cutoff $-0.5$ yr & $+0.267$ & $0.122$ & --- & 15/4; \emph{exceeds true cutoff} \\
cutoff $+0.5$ yr & $+0.046$ & $0.070$ & --- & 19/4 \\
cutoff $+1.0$ yr & $-0.311$ & $0.074$ & --- & 10/3 \\
\midrule
\multicolumn{5}{l}{\emph{Panel D: automated \natex{} runs}}\\
SuDDDS DD (94-cell panel) & $+0.232$ & $0.088$ & refused & 0 in-space placebos; gate closed \\
DiK slope, cutoff 0, bw 0.75 & $+0.015$ & $0.443$ & --- & drifts to $-0.843$ (0.218) at bw 2.0 \\
DiK placebo $+0.5$, bw 0.75 & $-1.467$ & $0.644$ & --- & exceeds all true-cutoff estimates \\
\bottomrule
\end{tabular}

\medskip
\parbox{0.92\textwidth}{\footnotesize Notes: jackknife over benchmarks keeps
$\tau \in [+0.085, +0.171]$ (positive in all 19 leave-one-out samples). RI
permutation quantiles (5\%/50\%/95\%): $-0.224$ / $+0.028$ / $+0.200$. The
SuDDDS scan $p=0.010$ localizes the declared treatment path and is
mechanical given the declared $\theta$; its refusal of a placebo $p$ is the
design working as intended with one binary group dimension.}
\end{table}

\section{Caveats and Conclusion}\label{sec:conclusion}

Six caveats bound the claims, none softened. First, \emph{four held-out
clusters}: CR1/$t(18)$ inference alone is not trustworthy at this cluster
count \cite{cameron2015}; the wild-cluster bootstrap and the exact
label-randomization test are the inference of record, and they do not
reject --- but randomization inference assumes label exchangeability, which
the composition differences below strain; it remains the most honest
finite-sample check available here. Second, \emph{difficulty confounding}:
held-out status is not randomly assigned --- the held-out set is the
hardest-benchmark set, and post~$\times$~difficulty controls absorb most of
the premium. Third, \emph{mechanical attenuation}: ECI is trained on the
public benchmarks, so contamination inflates the capability control itself;
the ability-proxy variants quantify this (tripling $\tau$) but inherit every
other weakness. Fourth, \emph{localization failure}: the $-0.5$-year
placebo contrast exceeds the true-cutoff contrast, so the design cannot
distinguish a release-localized step from a smooth group trend difference
--- and reads the data as the latter. Fifth, \emph{time confounding}:
within a benchmark, event time and calendar time are collinear; the
release-year-FE specification leans entirely on cross-benchmark stagger.
Sixth, \emph{declared inputs and ceilings}: three hand-filled release dates
and the held-out membership are declared, not estimated; saturation
ceilings compress public post-release residuals \emph{against} the
hypothesis, and the logit and drop-saturated variants only partially derate
this, so a true premium could be somewhat attenuated --- but ceilings cannot
manufacture the placebo-cutoff pattern.

The verdict is a suggestive-sign null. Every level specification is
positive; the held-out post-release shortfall ($-0.19$~z) is exactly the
shape leakage or test-task training \cite{dominguez2024} would leave; and
the ability-proxy estimates show the premium is larger once the
contaminated capability index is removed from the control set. But the
premium is not significant under any cluster-honest test, collapses when
held-out difficulty is controlled, and --- decisively --- is not localized
at benchmark release. What these data support is at most a slow divergence
of public-benchmark performance from held-out performance across model
generations, consistent with gradual test-task adaptation
\cite{dominguez2024} and with \cite{zhang2024}'s finding that frontier
models show little sharp overfitting; what they do not support is a
detectable step premium that switches on when a benchmark's answers become
available. Benchmark designers' held-out and rotating designs
\cite{glazer2024,phan2025,chollet2025,white2024} are the reason this test
was possible at all; more held-out clusters would make it decisive.

\paragraph{Reproducibility.} All estimates derive from the public Epoch AI
tables plus three declared release dates; no dataset files are committed.
The only stochastic step is the wild-cluster bootstrap (single generator,
seed 0); randomization inference is exhaustive and \natex{} runs declared
seed 0. Figure~\ref{fig:fig1} regenerates from the committed
\texttt{figures/make\_fig.py}, which asserts the headline estimates, the
placebo grid, the jackknife range, and the panel composition against the
numbers of record before drawing.

\footnotesize
\begin{thebibliography}{9}
\setlength{\itemsep}{1pt}\setlength{\parskip}{0pt}

\bibitem{bockerman2025}
B\"ockerman, P., Jysm\"a, S., and Kanninen, O. (2025).
\emph{Difference-in-Kinks Design}.
IZA Discussion Paper No.~18313.
\url{https://docs.iza.org/dp18313.pdf}

\bibitem{cameron2015}
Cameron, A.~C., and Miller, D.~L. (2015).
A practitioner's guide to cluster-robust inference.
\emph{Journal of Human Resources}, 50(2), 317--372.

\bibitem{chollet2025}
Chollet, F., Knoop, M., Kamradt, G., and Landers, B. (2025).
ARC-AGI-2: A new challenge for frontier AI reasoning systems.
arXiv:2505.11831.

\bibitem{dominguez2024}
Dominguez-Olmedo, R., Dorner, F.~E., and Hardt, M. (2024).
Training on the test task confounds evaluation and emergence.
arXiv:2407.07890.

\bibitem{epochai}
Epoch AI (2026).
Epoch Capabilities Index (ECI) and Benchmarking Hub data.
\url{https://epoch.ai/benchmarks} (panel extracted June 2026).

\bibitem{glazer2024}
Glazer, E., Erdil, E., Besiroglu, T., et al. (2024).
FrontierMath: A benchmark for evaluating advanced mathematical reasoning in
AI.
arXiv:2411.04872.

\bibitem{natex}
Hillebrandt, H. (2026).
\emph{natex: automated natural-experiment discovery and estimation} (version
0.2.0). Software.
\url{https://github.com/HaukeHillebrandt/natex}

\bibitem{ho2025}
Ho, A., Denain, J.-S., Atanasov, D., Albanie, S., and Shah, R. (2025).
A Rosetta Stone for AI benchmarks.
arXiv:2512.00193.

\bibitem{oren2023}
Oren, Y., Meister, N., Chatterji, N., Ladhak, F., and Hashimoto, T.~B.
(2023).
Proving test set contamination in black box language models.
arXiv:2310.17623.

\bibitem{phan2025}
Phan, L., et al. (2025).
Humanity's Last Exam.
arXiv:2501.14249.

\bibitem{roberts2023}
Roberts, M., Thakur, H., Herlihy, C., White, C., and Dooley, S. (2023).
Data contamination through the lens of time.
arXiv:2310.10628.

\bibitem{white2024}
White, C., Dooley, S., Roberts, M., et al. (2024).
LiveBench: A challenging, contamination-limited LLM benchmark.
arXiv:2406.19314.

\bibitem{zhang2024}
Zhang, H., Da, J., Lee, D., et al. (2024).
A careful examination of large language model performance on grade school
arithmetic.
arXiv:2405.00332.

\end{thebibliography}

\end{document}
